\section{Chapter~\ref{sec:neurontypes}. Neuron Types}

The chapter's numbers rest on direct cell counts in the human brain where such counts exist; the ``one neuron, one transmitter'' rule rests on cell atlases and on experiments that have also found the exceptions; the role of \nt{DOPA} rests on spike recordings; and the roles of the other broadcast channels rest on hypotheses.

{\footnotesize
\setlength{\tabcolsep}{3pt}\renewcommand{\arraystretch}{1.15}
\begin{longtable}{>{\raggedright}p{0.5cm}>{\raggedright}p{4.0cm}>{\raggedright}p{5.6cm}>{\raggedright}p{2.3cm}>{\raggedright\arraybackslash}p{1.7cm}}
\toprule
No. & Claim & Experiment and what was measured & Source & Status \\
\midrule\endfirsthead
\toprule
No. & Claim & Experiment and what was measured & Source & Status \\
\midrule\endhead
1 & The human brain has about $8.6\cdot10^{10}$ neurons, and almost all \nt{GLUT} neurons are small cerebellar neurons (Table~\ref{tab:types}) & Isotropic fractionator: tissue is dissolved into a uniform suspension of nuclei, and nuclei carrying the neuronal marker NeuN are counted. Adult men have $86.1\pm8.1$ billion neurons; only $19\,\%$ of them are in the cerebral cortex, and the bulk is in the cerebellum & \cite{Azevedo2009} & measured \\
2 & Almost every neuron has one principal transmitter (Proposition~\ref{prop:dale}), but there are exceptions & Transcriptomic atlas of the mouse brain: about $4\cdot10^6$ cells, $5322$ clusters; each cluster is assigned a transmitter from its synthesis and packaging genes. The exceptions were measured directly: in slices, \nt{DOPA} neurons also release GABA through the same vesicular transporter and inhibit striatal neurons; in some \nt{DOPA} axons, glutamate and dopamine leave from different terminals of the same wire (electron microscopy and optogenetics) & \cite{Strata1999,Yao2023,Tritsch2012,Zhang2015} & measured \\
3 & The entire column of the weight matrix has one sign~\eqref{eq:dalesign} & Derived in the chapter from Proposition~\ref{prop:dale} and the fact that glutamate receptors are almost all excitatory and GABA receptors inhibitory. There is no direct test of ``all recipients of one neuron receive a weight of one sign'' as a general rule; counterexamples are GABA co-release by \nt{DOPA} neurons (row~2) and recipients with receptors of opposite sign (row~11) & derivation in the chapter & theory \\
4 & Three-quarters to four-fifths of cortical neurons are \nt{GLUT}, a fifth to a quarter are \nt{GABA} & GABA immunostaining in columns $50\um$ wide through the full depth of the cortex in $10$ areas of monkey cortex: in $8$ areas GABA neurons make up about $25\,\%$ of all neurons, in primary visual cortex less than $20\,\%$ & \cite{Hendry1987} & measured \\
5 & A \nt{GLUT} neuron has about $10^4$ outputs & Postmortem stereology: the neocortex of young men has $1.64\cdot10^{14}$ synapses (electron microscopy, disector) and about $2.3\cdot10^{10}$ neurons. The ratio gives $\sim7\cdot10^3$ synapses per neuron; on average the number of inputs equals the number of outputs & \cite{Tang2001synapses,Pakkenberg1997} & measured \\
6 & The output of a \nt{DOPA} neuron branches into $10^5$--$10^6$ terminals; this is an estimate from target coverage, not a count & Viral membrane labeling of single rat \nt{DOPA} neurons and full axon reconstruction: one neuron covers $0.45$ to $5.7\,\%$ (on average $2.7\,\%$) of the striatal volume. Coverage was measured, not the number of terminals: that is inferred from coverage to order of magnitude, with no direct terminal count. For \nt{SERO} and \nt{NORA} the terminal numbers are rough estimates & \cite{Matsuda2009} & measured \\
7 & Broadcast neurons are few: \nt{DOPA} $\sim5\cdot10^5$ (the main one of two groups, a lower bound), \nt{SERO} $\sim3\cdot10^5$ (sum of two partial counts), \nt{HIST} $\sim6\cdot10^4$ (Table~\ref{tab:types}, Fig.~\ref{fig:typepie}) & Stereological counts in humans: $550\,000$ pigmented neurons in the substantia nigra (without the ventral tegmental area); $235\,000$ neurons in the dorsal raphe nucleus (all neurons of the nucleus) and another $125\,000$ serotonin neurons in the pontine tegmentum; about $64\,000$ histamine neurons in the hypothalamus. For \nt{NORA}, \nt{GLYC}, \nt{ACET}, and \nt{ADRE} there is no direct count: in the table these are rough order-of-magnitude estimates & \cite{Pakkenberg1991,Baker1990,Baker1991,Airaksinen1991} & measured \\
8 & Glutamate in the gap shoots up to about a millimolar and decays within $\sim1\ms$ & From the displacement of a rapidly unbinding competitive blocker from NMDA receptors during synaptic transmission (cultured hippocampal neurons): peak $1.1$\,mM, decay time constant $1.2\ms$ & \cite{Clements1992} & measured \\
9 & In working cortex the excitatory and inhibitory currents are nearly balanced, and the neuron sits near threshold & Theory: a network with strong connections settles into the balanced regime on its own. Experiment: in ferret prefrontal cortex in vivo, during ``up'' states the reversal potential of the synaptic current stays near $-37$\,mV for hundreds of milliseconds, although the conductance changes by $\sim21$\,nS: excitation and inhibition rise and fall together & \cite{vanVreeswijk1996,Haider2006} & measured \\
10 & \nt{DOPA} neurons report the reward-prediction error~\eqref{eq:rpe} (Fig.~\ref{fig:rpe}) & Spikes of monkey \nt{DOPA} neurons: a burst for unexpected juice; after learning, a burst for the cue and no response to predicted juice; a dip in rate when promised juice is withheld. In a quantitative experiment, the rate is proportional to the difference between the current reward and a weighted average of past ones, but only for positive errors. Equation~\eqref{eq:rpe} itself is the temporal-difference algorithm & \cite{Schultz1997,Bayer2005,SuttonBarto2018} & measured \\
11 & The receptor sets the sign and speed: ionotropic, fractions of a millisecond; metabotropic, much slower (Proposition~\ref{prop:receptor}) & Brief pulses of glutamate on membrane patches from hippocampal neurons: the current through ionotropic receptors rises (from $20$ to $80\,\%$) in $0.2$--$0.6\ms$ and decays in $2$--$3\ms$. The slow inhibitory potential of pyramidal neurons is abolished by a selective blocker of the metabotropic GABA receptor. Two families of dopamine receptors with opposite action; in the striatum, neurons with the D1 receptor need dopamine to strengthen synapses, neurons with D2 to weaken them & \cite{Colquhoun1992,DutarNicoll1988,Beaulieu2011,Shen2008} & measured \\
12 & A broadcast channel is not one wire but a bundle of a few lines (Section~\ref{sec:buses}) & Viral-genetic labeling of serotonin neurons in the mouse dorsal raphe nucleus: neurons projecting to the amygdala and to the frontal cortex lie in different parts of the nucleus, branch into complementary regions, and respond oppositely to an aversive stimulus. Review: subgroups of locus coeruleus neurons (\nt{NORA}) encode different processes & \cite{Ren2018,Poe2020} & measured \\
13 & \nt{NORA} sets the gain $g$ & The adaptive gain hypothesis; a network model in which catecholamines increase the steepness of the transfer characteristic. There is no direct measurement of ``$g$ rises with noradrenaline'' across the whole network; what has been measured is that pupil diameter follows the spikes of locus coeruleus neurons in the monkey & \cite{AstonJones2005,ServanSchreiber1990,Joshi2016} & hypothesis \\
14 & \nt{ACET} switches ``write/read'' and sets the learning rate & Hypothesis. Support: in rat olfactory cortex slices, acetylcholine agonists ($100$\,\textmu M) weaken the internal, ``recalling'' synapses by more than $60\,\%$ and the input synapses by less than $15\,\%$ & \cite{Hasselmo2006,HasselmoBower1992} & hypothesis \\
15 & \nt{SERO} sets the discount factor $\gamma$; serotonin neurons fire steadily, a few spikes per second, and are almost silent in one of the phases of sleep & The ``serotonin $=\gamma$'' hypothesis. Strongest argument: optogenetic activation of serotonin neurons in the mouse dorsal raphe nucleus reduces the number of premature give-ups while waiting for a delayed reward. The firing rate and silence during rapid-eye-movement sleep have been measured (review of recordings) & \cite{Doya2002,Miyazaki2014,JacobsAzmitia1992} & hypothesis \\
\bottomrule
\end{longtable}}
